\begin{lemma}\label{lem:octic-dyadic-F4-kernel}
Let $A/\mathbb Z_2$ be an abelian scheme with commuting
endomorphisms $T$ and $\delta$ satisfying
$T_{\mathbb F_2}=F_2+\delta V_2$. Let $R$ be the ring of
integers of the unramified quadratic extension of $\mathbb Q_2$.
Then $T^2-\delta^2$ kills the torsion kernel of reduction
$A(R)\to A(\mathbb F_4)$.
\end{lemma}
\begin{proof}
The torsion kernel has exponent $2$. The schematic closure
of a nonzero point and zero is finite flat of rank $2$,
with connected special fibre. The rank-two Tate--Oort
classification over the unramified ring $R$ identifies it
with $\mu_2$. Reduction on
$\operatorname{Hom}_R(\mu_2,A[2])$ is injective: after
Cartier duality, maps to the constant \'etale group $\mathbb Z/2$
are determined by their open and closed fibres, hence by
their reductions over the henselian ring $R$.

Since $F_2V_2=[2]$, we have the exact identity
\[
 T^2=F_4+[4]\delta+\delta^2V_4.
\]
On a copy of $\mu_2$ defined over $\mathbb F_4$, $F_4$ is
zero, $V_4$ is the identity, and $[4]$ is zero.
Naturality therefore gives
$(T^2-\delta^2)\overline\iota=0$ for the inclusion
$\iota:\mu_2\hookrightarrow A[2]$. Injectivity of reduction
on this Hom group proves the assertion over $R$.
\end{proof}

\begin{lemma}\label{lem:octic-F4-Cartier}
Let $F=\mathbb Q(\zeta_3)$ and let $Y/\mathbb Q$ be a
modular curve with good reduction at $2$. Suppose that
$\operatorname{rk}\operatorname{Jac}(Y)(F)=0$, that $Y_F$
has an $F$-rational cusp, and that its cusps split over a
cyclotomic extension unramified at $2$. Suppose that $T_2$
and $\delta_2$ descend to $\operatorname{Jac}(Y)$ and
satisfy Eichler--Shimura. Put
\[
 U=(\mathcal C^2-1)H^0(Y_{\mathbb F_4},\Omega^1),
\]
where $\mathcal C$ is Cartier. If restriction from $U$
to every effective cuspidal degree-$d$ divisor over
$\mathbb F_4$ is surjective, then no noncuspidal
effective $F$-rational degree-$d$ divisor has entirely
cuspidal reduction at $2$.
\end{lemma}
\begin{proof}
Let $D$ be such a divisor and let $C$ be the unique \'etale
cuspidal lift of its reduction. Although $C$ need not be
globally defined over $F$, its multiplicities are constant
on Frobenius-$4$ cycles. Thus $C$ is defined over the local
field $F_v$, the unramified quadratic extension of
$\mathbb Q_2$. It is also globally defined over a
decomposition field in the cyclotomic cusp field.

By rank zero and Manin--Drinfeld, $[D-C]$ is torsion.
It lies in the reduction kernel over $F_v$.
Lemma~\ref{lem:octic-dyadic-F4-kernel} gives
\[
 (T_2^2-\delta_4)[D-C]=0,
 \qquad \delta_4=\delta_2^2.
\]
On cotangents in characteristic two, the image of
$T_2^2-\delta_4$ is
$\delta_4^*(\mathcal C^2-1)H^0(\Omega^1)=U$.
Here $F_4$ pulls back to zero, $V_4$ pulls back to
$\mathcal C^2$, and the diamond operator commutes with
Cartier. The restriction hypothesis is therefore formal
immersion at the reduction of $C$ for
\[
 Y_F^{(d)}\longrightarrow\operatorname{Jac}(Y)_F,
 \quad E\longmapsto(T_2^2-\delta_4)[E-dc_0].
\]
The sections $D,C$ have the same reduction and image,
so formal immersion gives $D=C$, a contradiction.
\end{proof}

For computations over $\mathbb F_4$, Cartier is semilinear.
If its row matrix is $M$ in a basis of differentials, then
the row matrix of $\mathcal C^2$ is $M^{[2]}M$, where
$M^{[2]}$ denotes entrywise squaring. The program
\texttt{cartier2\_3\_15\_D4.m} verifies this identity
against direct double application of Cartier. Its
restriction computation has $616$ exceptions among
$5000$ divisors, so this lemma alone does not exclude
$C_3\times C_{15}$ over octic fields.
